\BeleskeID{09_3-s053}
% element: 09_3-s053:e04
% slide-id: 09_3-s053
Za I funkciju, kada je $A=1$, gornja grana prosleđuje $B$. Kada je $A=0$, donja grana prosleđuje nulu. Umesto stalne nule može se dovesti $A$, jer je u trenutku provođenja te grane $A$ upravo nula.

Za ILI funkciju, kada je $A=1$, gornja grana prosleđuje jedinicu. Kada je $A=0$, donja prosleđuje $B$. Umesto stalne jedinice može se koristiti $A$, jer gornja grana vodi samo kada je $A$ jedan. Prikazane jednakosti porede ove dve realizacije za svaku funkciju.
